\paragraph*{Infant Drug Testing by Race}

Across all years, 16.5\% of Black infants are tested compared with 7.3\% of White infants. The left panel of Figure \ref{fig:figure_1_testing_positive_by_year} shows that the testing gap was quite stable from 2019 to 2023. Despite the large gap in testing rates, the right panel of Figure \ref{fig:figure_1_testing_positive_by_year} shows that the positive testing rate was nearly identical among White and Black infants who were tested. The average positive testing rate for White infants was 19.7\% compared to 20.2\% among Black infants. The positive testing rate fell for both groups over time, particularly in 2022 and 2023. 

\paragraph*{Clinical Risk Differences}

Table \ref{tab:summary_table} reports demographic and clinical characteristics for the mother and infant by race. Black infants were more likely to be born premature and to have low birth weight than White infants. Mothers of Black infants were about 1 year younger and received more prenatal care than mothers of White infants. Nearly 59\% of the mothers of Black infants were covered by Medicaid compared with 26\% of White mothers. About 12\% of mothers of Black infants had a history of substance use disorder (SUD) and 6\% had a history of a chronic pain disorder (CPD); among mothers of White infants, 10\% had a previous SUD and 4\% had a previous CPD diagnosis. In contrast, the mothers of Black infants had lower rates of mental health disorder than the mothers of White infants. The Black infants in our sample reside in counties with higher drug overdose mortality than White infants (2.7 vs 2.03 overdose deaths per 10,000). However, the county rates of neonatal abstinence syndrome (NAS) were about the same for both White and Black infants.

\paragraph*{Testing Rates in Risk Sub-populations}

Figure \ref{fig:figure_testing_sud} presents infant drug testing rates by race and maternal SUD history. For both White and Black infants, drug testing rates are about twice as high among those with maternal SUD history. In this high risk group, the racial gap in testing rates has declined over time and is now quite small. Testing rates are lower among infants with no maternal SUD history, but the racial gap is substantial and more stable over time. Figure \ref{fig:figure_positive_sud} illustrates that positive testing rates are about the same for White and Black infants, regardless of maternal SUD history.

The relationship between infant drug testing rates and neonatal abstinence syndrome (NAS) rates at the county level is depicted in Figure \ref{fig:figure_nas_test}. Testing rates tend to be higher for both groups in counties that have higher rates of NAS, but the association is stronger for Black infants. Appendix Figure \ref{fig:figure_nas_positive} shows that the positivity rate is essentially flat with respect to the NAS rate across counties. Appendix Figures \ref{fig:figure_county_mort_test}
and \ref{fig:figure_county_mort_positive} show similar relationships with county drug mortality rates.

\paragraph*{Decomposing Racial Testing Gaps}

We used inverse propensity score weights to reweight the White and Black infant samples to match the covariate distribution of the other group. Appendix Table \ref{propensity_model} reports the logistic regression we used to estimate propensity scores, and Appendix Table \ref{tab:balance_reweight} shows covariate balance in the raw and reweighted sample.

Panel A of Table \ref{tab:gap_testing_ipw} shows that the observed testing rate was 7.3\% for White infants and 16.5\% for Black infants. The second row shows implied testing rates when the data are reweighted. When the sample of Black infants is reweighted to match the covariate distribution of the White infants, the testing rate in the matched Black sample falls to 12.9\%. When the White infants are reweighted to match the Black sample covariate distribution, the testing rate among White infants rises to 10.2\%.  

Figure \ref{fig:testing_rates_reweighted} shows how testing rates vary over time among White infants, Black infants, and Black infants reweighted to match the White covariate distribution. Setting clinically relevant individual and community risk factors among Black infants equal to their levels in the White infant sample substantially narrows the testing gap but does not eliminate it. 

Panel B of Table \ref{tab:gap_testing_ipw} decomposes the racial testing gap into the share explained by observed risk factors and the share that is unexplained. The observed testing gap is $\Delta_{B,W} = 16.5\% - 7.3\%$, which is $9.2$ percentage points. Reweighting the Black infants so that they have the same clinical risk distribution as the White sample but testing each group according to its own race-specific decision functions gives a structural gap of $\delta_s = 12.9\% - 7.3\% = 5.6$ percentage points. Applying the Black decision function to both White and Black infants leads to a risk-factor explained gap of $\delta_x = 16.5\% - 12.9\% = 3.5$ percentage points. This implies that clinical covariates account for about $39\%$ of the raw testing gap. The remaining $61\%$ of the gap may reflect either risk factors that are observed by health care workers but are not observed in the EMR data, or structural differences in the way that health care workers interpret the same characteristics differently for White and Black mothers and infants. \footnote{Quantifying the share of the testing gap attributable to observable risks and unmeasured factors may be sensitive to which group is used as the reference category. In this case, the results are not very sensitive to the reference group: Panel C of Table \ref{tab:gap_testing_ipw} shows that observed risk factors explain 33\% of the gap when we change the reference group.}